% !TeX root = ../document_maitre.tex

\label{juridique}

Les enjeux techniques de l'\gls{IA} ne repésente qu'une fraction des considérations à prendre en compte. Leurs capacités d'entraînement et de traitement sur des milliers de données variées, ont entraînés une très vive popularité de ces outils. 

Cette dimension, mêlée à leur opacité amène à définir un cadre législatif international, en plein essor. Ce dernier fusionne avec le cadre juridique patrimonial afin de normer et contraindre les outils que nous envisageons dans leur conception et leur traitement.

\subsection{Cadre juridique du patrimoine et des données patrimoniales}

Le patrimoine est basé sur des documents et objets comprenant des informations au sens large, issues d'institutions publiques ou du cadre privé. L'enjeu réside dans le sens de ces informations, pouvant être de nature sensible. Ce caractère particulier \guillemet{s'oppose} à l'obligation de communication du patrimoine. 

Ainsi, les textes juridiques cherchent à équilibrer cette diffusion du patrimoine avec la protection des informations sensibles. En France, le \textit{Code du Patrimoine}\footcite{CodePatrimoineLegifrance} est la pierre angulaire du cadre juridique. Il impose un délai de communication des documents en fonctions de leur typologie et contenu.\footcite[Article 213.1 à 213.3]{CodePatrimoineLegifrance}Et régit également les conditions de dépôt et de communication des archives et patrimoines privés, presque sur mesure en fonction des cas.\footcite[Article 213-6]{CodePatrimoineLegifrance}

Ce code est complété par le \textit{Code de la propriété intellectuelle} avec notamment le droit d'auteur.\footcite{DroitDauteurINPI2024} Le créateur d'une \oe{uvre} possède dès lors des droits moraux\footcite[Article 121-1 à 121-9]{DroitDauteurINPI2024}, protégeant le lien entre l'\oe{uvre} et son créateur, et des droits patrimoniaux, disons d'exploitation de l'\oe{uvre}\footcite[Article 122-1 à 122-12]{DroitDauteurINPI2024}.\newline

L'informatisation des collections patrimoniales a amené à compléter l'appareil juridique patrimonial avec un volet spécialisé sur les données informatiques. On peut penser au \textit{Code des relations entre le public et l'administration}. Il encadre la réutilisation des documents informatisés à des fins de diffusion, de traitement algorithmique et d'entraînement de \glspl{mod} d'\gls{IA}.\footcite{LivreIIILACCES}

Mais surtout au texte européen : \textit{Règlement général sur la protection des données}\footcite{ReglementGeneralProtection}, ou RGPD. Il est particulièrement intéressant dans notre position. Il définit les obligations qui incombent aux entreprises comme SkinSoft qui stockent des données dont elles ne sont pas propriétaires. Les données sont classées par catégories de sensibilités auxquelles sont associés des traitements à effectuer avant diffusions ou conservation. Cette dernière est également contrôlée dans sa forme et sa durée.

Ainsi, les fournisseurs de \gls{sgc} ont la position de \guillemet{sous-traitant}.\footcite[Article 28]{ReglementGeneralProtection} L'entreprise n'a pas de droits sur les données conservées et le propriétaire dispose d'un accès permanents à ces dernières. Tout traitement à effectuer doit être approuvé par l'institution ou la personne propriétaire des données.

\subsection{Contraintes juridiques de l'\gls{IA}}

Le cadre juridique se dessine avec le texte fondateur promulgué par l'Union européenne en 2024 : \textit{l'AI Act}\footcite{LoiEuropeenneLintelligence}. Il impose des principes généraux sur ces outils. À savoir : la transparence des jeux d'entraînement, des traitements et du stockage des données ou encore le droit d'opposition accordé aux propriétaires des données. L'objectif est de sécuriser l'utilisation d'outils d'\gls{IA} sur les données. En fonction de la sensibilité des données et du traitement appliqué, un niveau de risque de l'outil est déterminé. Ce dernier impose des principes et conditions de sécurité.

La \guillemet{recherche augmentée}, le \gls{chatbot} et l'indexation semi-automatisée relèvent du risque limité. Cela impose une transparence de l'entreprise dans l'utilisation, le traitement et le stockage des données des institutions tout en respectant le droit d'opposition. 

Cela s'ajoute au respect du \textit{Code du patrimoine} et du \textit{RGPD} précédemment cités. Ainsi, les outils envisagés doivent : respecter les délais de communications, traiter les données sensibles dans l'anonymat, ne pas diffuser les données sans l'accord des propriétaires et être transparent sur les traitements ainsi que la conservation des données. Tout cela doit être encadré par une infrastructure sécurisée. \newline

Ce cadre juridique est hautement évolutif. Aujourd'hui, nous ne sommes qu'à un stade précoce de construction de ce cadre. L'\textit{AI Act} instaure les principes de fond à mettre en relation avec l'appareil juridique patrimonial. Sur ces principes fondamentaux s'éleveront certainement de nouvelles législations. Ainsi, la conformité juridique en matière d'\gls{IA} est une question évolutive. Et il incombe à des structures comme SkinSoft, si le projet de création d'outils d'\gls{IA} aboutit, d'évoluer avec ce cadre. 